\chapter{Formulae}

\small
\subsection{Combinatorics}

\begin{enumerate}
    \item $\displaystyle \displaystyle \sum \limits_{0\leq k \leq n} {n-k \choose k} = Fib_{n+1}$
    \item $\displaystyle \displaystyle {n \choose k}={n \choose n-k}$
    \item $\displaystyle \displaystyle {n \choose k}+{n \choose k+1}={n+1 \choose k+1}$
    \item $\displaystyle \displaystyle k{n \choose k}=n{n-1 \choose k-1}$
    \item $\displaystyle \displaystyle {n \choose k}=\frac{n}{k}{n-1 \choose k-1}$
    \item $\displaystyle \sum \limits_{i=0}^n{n \choose i}=2^n$
    \item $\displaystyle \sum \limits_{i\geq 0}{n \choose 2i}=2^{n-1}$
    \item $\displaystyle \sum \limits_{i\geq 0}{n \choose 2i+1}=2^{n-1}$
    \item $\displaystyle \sum \limits_{i= 0}^k \left( -1 \right) ^i{n \choose i}=\left( -1 \right) ^k{n-1 \choose k}$
    \item $\displaystyle \sum \limits_{i= 0}^k{n+i \choose i}= \sum \limits_{i= 0}^k{n+i \choose n} = {n+k+1 \choose k}$
    \item $\displaystyle \displaystyle1{n \choose 1}+2{n \choose 2}+3{n \choose 3}+\dots+n{n \choose n}=n2^{n-1}$
    \item $\displaystyle \displaystyle1^2{n \choose 1}+2^2{n \choose 2}+3^2{n \choose 3}+\dots+n^2{n \choose n}=(n+n^2)2^{n-2}$
    \item \textbf{Vandermonde’s Identify:} \\ $\displaystyle \sum \limits_{k=0}^r{m \choose k}{n \choose r-k}={m+n \choose r}$
    \item \textbf{Hockey-Stick Identify:} \\ $\displaystyle \ n,r \in N, n > r, \sum \limits_{i=r}^n{i \choose r}={n+1 \choose r+1}$
    \item $\displaystyle \sum \limits_{i=0}^k{k \choose i}^2={2k \choose k}$
    \item $\displaystyle \sum \limits_{k=0}^n{n \choose k}{n \choose n-k}={2n \choose n}$
    \item $\displaystyle \sum \limits_{k=q}^n{n \choose k}{k \choose q}=2^{n-q}{n \choose q}$
    \item $\displaystyle \sum \limits_{i=0}^nk^i{n \choose i}=(k+1)^n$
    \item $\displaystyle \sum \limits_{i=0}^n{2n \choose i}=2^{2n-1}+\frac{1}{2}{2n \choose n}$
    \item $\displaystyle \sum \limits_{i=1}^n{n \choose i}{n-1 \choose i-1}={2n-1 \choose n-1}$
    \item $\displaystyle \sum \limits_{i=0}^n{2n \choose i}^2=\frac{1}{2} \left( {4n \choose 2n}+{2n \choose n}^2 \right)$
    \item \textbf{Highest Power of $\displaystyle 2$ that divides $^{2n}C_n$:}
    Let $\displaystyle x$ be the number of $\displaystyle 1$s in the binary representation. Then the number of odd terms will be $\displaystyle 2^x$. Let it form a sequence. The $n$-th value in the sequence (starting from $n$ = 0) gives the highest power of $\displaystyle 2$ that divides $^{2n}C_n$.
    \item \textbf{Pascal Triangle}
    \begin{enumerate}
        \item In a row $\displaystyle p$ where $\displaystyle p$ is a prime number, all the terms in that row except the $\displaystyle 1$s are multiples of $\displaystyle p$.
        \item Parity: To count odd terms in row $n$, convert $n$ to binary. Let $\displaystyle x$ be the number of $\displaystyle 1$s in the binary representation. Then the number of odd terms will be $\displaystyle 2^x$.
        \item Every entry in row $\displaystyle 2^n-1, n\geq 0,$ is odd.
    \end{enumerate}
    \item An integer $n\geq 2$ is prime if and only if all the intermediate binomial coefficients \\ \\  $\displaystyle \displaystyle {n \choose 1},{n \choose 2},\dots, {n \choose n-1}$ are divisible by $n$.
    \item \textbf{Kummer’s Theorem:} For given integers $n\geq m\geq 0$ and a prime number $\displaystyle p$, the largest power of $\displaystyle p$ dividing $\displaystyle \displaystyle {n \choose m}$ is equal to the number of carries when $m$ is added to $n$-$m$ in base $\displaystyle p$. For implementation take inspiration from lucas theorem.
    \item Number of different binary sequences of length $n$ such that no two $0$’s are adjacent=$\displaystyle Fib_{n+1}$
    \item \textbf{Combination with repetition:} Let’s say we choose $\displaystyle \displaystyle k$ elements from an $n$-element set, the order doesn’t matter and each element can be chosen more than once. In that case, the number of different combinations is: $\displaystyle \displaystyle {n+k-1 \choose k}$
    \item Number of ways to divide $n$ persons in $\displaystyle \frac{n}{k}$ equal groups i.e. each having size $\displaystyle \displaystyle k$ is
    \[\displaystyle \frac{n!}{k!^{\frac{n}{k}} \left( \frac{n}{k} \right) !}= \prod \limits_{n \geq k}^{n-=k}{n-1 \choose k-1}\]
    \item The number non-negative solution of the equation: $\displaystyle x_1+x_2+x_3+\dots+x_k=n \text{ is}{n+k-1 \choose n}$
    \item Number of ways to choose $n$ ids from $\displaystyle 1$ to b such that every id has distance at least k $\displaystyle \displaystyle =\left( \frac{b-(n-1)(k-1)}{n} \right)$
    \item $\displaystyle \displaystyle \sum \limits_{i=1,3,5,\dots}^{i\leq n} \binom{n}{i} a^{n-i}b^i = \frac{1}{2} ((a+b)^n-(a-b)^n)$
    \item $\displaystyle \displaystyle \sum \limits_{i=0}^{n} \dfrac{\binom{k}{i}}{\binom{n}{i}} = \dfrac{\binom{n+1}{n-k+1}}{\binom{n}{k}}$
    \item \textbf{Derangement:} a permutation of the elements of a set, such that no element appears in its original position. Let $d(n)$ be the number of derangements of the identity permutation fo size $n$.
    \[\begin{array}{@{}l}d(n) = (n - 1)\bigl(d(n - 1) + d(n - 2)\bigr),\\ \text{where } d(0)=1,\ d(1)=0.\end{array}\]
    \item \textbf{Involutions:} permutations such that $\displaystyle p^2 =$ identity permutation. $\displaystyle a_0 = a_1 = 1$ and $\displaystyle a_n=a_{n-1} + (n-1)a_{n-2}$ for $n>1$.
    \item Let $T(n,k)$ be the number of permutations of size $n$ for which all cycles have length $\displaystyle \leq k$.
    \[T(n, k)=
    \begin{cases}
      n! & ; n \le k\\
      n \cdot T(n-1, k) \\ - F(n-1, k) \\ \cdot T(n-k-1, k) & ;n > k \\
    \end{cases}\]
    Here $\displaystyle F(n, k) = n \cdot (n - 1) \cdot \dots \cdot (n - k + 1)$
    \item \textbf{Lucas Theorem}
    \begin{enumerate}
      \item If $\displaystyle p$ is prime, then $\displaystyle \left(\frac{p^{a}}{k}\right) \equiv 0 (\mod\; p)$
      \item For non-negative integers $m$ and $n$ and a prime $\displaystyle p$, the following congruence relation holds:
      \(\displaystyle \left(\frac{m}{n}\right) \equiv \prod_{i=0}^{k} \left(\frac{m_{i}}{n_{i}}\right) (\text{mod}\; p),\)
      where,
      $m=m_{k} p^{k} +m_{k-1} p^{k-1}+\dots+m_{1} p+m_{0}$,
      and
      $n=n_{k}p^{k}+n_{k-1}p^{k-1}+\dots+n_{1}p+n_{0}$
      are the base $\displaystyle p$ expansions of $m$ and $n$ respectively. This uses the convention that $\displaystyle \left(\frac{m}{n}\right)=0$,when $m<n$.
    \end{enumerate}
    \item $\displaystyle \sum_{i=0}^{n}  {n \choose i } \cdot i^{k}$
    \\= $\displaystyle \sum_{i=0}^{n}  {n \choose i } \cdot \sum_{j=0}^{k}{ \left\{ \begin{matrix} k\cr j \end{matrix} \right\} \cdot i^{\underline{j}}}$
    \\= $\displaystyle \sum_{i=0}^{n} {n \choose i } \cdot \sum_{j=0}^{k} \left\{ \begin{matrix} k\cr j \end{matrix} \right\} \cdot  j! {n \choose i }$
    \\= $\displaystyle \sum_{i=0}^{n} \frac{n!}{(n-i)!} \cdot \sum_{j=0}^{k} \left\{ \begin{matrix}k\cr j \end{matrix} \right\} \cdot \frac{1}{(i-j)!}$
    \\= $\displaystyle \sum_{i=0}^{n} \sum_{j=0}^{k} \frac{n!}{(n-i)!} \cdot  \left\{\begin{matrix} k\cr j \end{matrix} \right\} \cdot \frac{1}{(i-j)!}$
    \\= n!$\displaystyle \sum_{i=0}^{n} \sum_{j=0}^{k} \left\{ \begin{matrix} k\cr j \end{matrix} \right\} \cdot \frac{1}{(n-i)!} \cdot \frac{1}{(i-j)!}$
    \\= n!$\displaystyle \sum_{i=0}^{n} \sum_{j=0}^{k} \left\{ \begin{matrix} k\cr j \end{matrix} \right\} \cdot { n-j \choose n-i} \cdot \frac{1}{(n-j)!}$
    \\= n!$\displaystyle \sum_{j=0}^{k} \left\{ \begin{matrix} k\cr j \end{matrix} \right\} \cdot \frac{1}{(n-j)!} \sum_{i=0}^{n} \cdot { n-j \choose n-i}$
    \\= $\displaystyle \sum_{j=0}^{k} \left\{\begin{matrix} k\cr j \end{matrix} \right\} \cdot n^{\underline{j}} \cdot 2^{n-j}$

    Here $n^{\underline{j}}=P(n,j)=\dfrac{n!}{(n-j)!}$ and $\displaystyle \left\{ \begin{matrix} k\cr j \end{matrix} \right\}$ is stirling number of the second kind.

    So, instead of $O(n)$, now you can calculate the original equation in $O(k^2)$ or even in $O(k \log^2 n)$ using NTT.
    \item $\displaystyle \displaystyle \sum \limits_{i=0}^{n-1} {i \choose j} x^i = x^j(1-x)^{-j-1} \\ \cdot \left( 1- x^n \sum \limits_{i=0}^j {n \choose i} x^{j-i} (1-x)^i \right)$
    \item $\displaystyle x_{0}, x_{1}, x_{2}, x_{3}, \dots , x_{n}$
    \\ $\displaystyle x_{0} + x_{1}, x_{1} + x_{2}, x_{2}+x_{3}, \dots x_{n}$
    \\ \dots
    \\ If we continuously do this $n$ times then the polynomial of the first column of the $n$-th row will be
    \[\displaystyle p(n)=\sum_{k=0}^{n}{n \choose k} \cdot x(k)\]
    \item If $\displaystyle P(n)=\sum_{k=0}^{n}{n \choose k} \cdot Q(k)$, then,
    \[Q(n)=\sum_{k=0}^{n}(-1)^{n-k}{n \choose k} \cdot P(k)\]
    \item If $\displaystyle P(n)=\sum_{k=0}^{n}(-1)^{k}{n \choose k} \cdot Q(k)$ , then,
    \[Q(n)=\sum_{k=0}^{n}(-1)^{k}{n \choose k} \cdot P(k)\]
\end{enumerate}

\subsection{Catalan Numbers}

\begin{enumerate}[resume]
    \item $\displaystyle C_n=\frac{1}{n+1}{2n \choose n}$
    \item $\displaystyle C_0=1,C_1=1\text{ and }C_n=\sum \limits_{k=0}^{n-1}C_k C_{n-1-k}$
    \item Number of correct bracket sequence consisting of $n$ opening and $n$ closing brackets.
    \item The number of ways to completely parenthesize $n$+$\displaystyle 1$ factors.
    \item The number of triangulations of a convex polygon with $n$+$\displaystyle 2$ sides (i.e. the number of partitions of polygon into disjoint triangles by using the diagonals).
    \item The number of ways to connect the $\displaystyle 2n$ points on a circle to form $n$ disjoint i.e. non-intersecting chords.
    \item The number of monotonic lattice paths from point $\displaystyle (0,0)$ to point $\displaystyle (n,n)$ in a square lattice of size $n\times n$, which do not pass above the main diagonal (i.e. connecting $\displaystyle (0,0)$ to $\displaystyle (n,n)$).
    \item The number of rooted full binary trees with $n$+$\displaystyle 1$ leaves (vertices are not numbered). A rooted binary tree is full if every vertex has either two children or no children.
    \item Number of permutations of $\{1, \dots, n\}$ that avoid the pattern $123$ (or any of the other patterns of length $3$); that is, the number of permutations with no three-term increasing subsequence. For $n = 3$, these permutations are $132, 213, 231, 312$ and $321$. For $n = 4$, they are $1432, 2143, 2413, 2431, 3142, 3214, 3241, 3412, 3421, 4132, 4213, 4231, 4312$ and $4321.$
    \item \textbf{Balanced Parentheses count with prefix:}
    The count of balanced parentheses sequences consisting of $n+k$ pairs of parentheses where the first $\displaystyle \displaystyle k$ symbols are open brackets. Let the number be $\displaystyle C_n^{(k)}$, then
    \[\displaystyle \displaystyle C_n^{(k)} = \frac{k+1}{n+k+1} \binom{2n+k}{n}\]
\end{enumerate}

\subsection{Narayana numbers}

\begin{enumerate}[resume]
    \item $N(n,k)=\frac{1}{n}\left(\frac{n}{k}\right)\left(\frac{n}{k-1}\right)$
    \item The number of expressions containing $n$ pairs of parentheses, which are correctly matched and which contain $\displaystyle \displaystyle k$ distinct nestings. For instance, $N(4, 2) = 6$ as with four pairs of parentheses six sequences can be created which each contain two times the sub-pattern ‘()’.
\end{enumerate}

\subsection{Stirling numbers of the first kind}

\begin{enumerate}[resume]
    \item The Stirling numbers of the first kind count permutations according to their number of cycles (counting fixed points as cycles of length one).
    \item $S(n,k)$ counts the number of permutations of $n$ elements with $\displaystyle \displaystyle k$ disjoint cycles.
    \item $S(n,k)=(n-1) \cdot S(n-1,k)+S(n-1,k-1),$
    \(where,\; S(0,0)=1,S(n,0)=S(0,n)=0\)
    \item $\displaystyle \displaystyle\sum_{k=0}^{n}S(n,k) = n!$
    \item The unsigned Stirling numbers may also be defined algebraically, as the coefficient of the rising factorial:
    \[\displaystyle x^{\bar{n}} = x(x+1)\dots(x+n-1) = \sum_{k=0}^{n}{ S(n, k) x^k}\]
    \item Lets $[n, k]$ be the stirling number of the first kind, then
    \[\displaystyle \bigl[\!\begin{smallmatrix} n \\ n\ -\ k \end{smallmatrix}\!\bigr] = \sum_{0 \leq i_1 < i_2 < i_k < n}{i_1i_2\dotsi_k.}\]
\end{enumerate}

\subsection{Stirling numbers of the second kind}

\begin{enumerate}[resume]
    \item Stirling number of the second kind is the number of ways to partition a set of n objects into k non-empty subsets.
    \item $S(n,k)=k \cdot S(n-1,k)+S(n-1,k-1)$,
    \(where \; S(0,0)=1,S(n,0)=S(0,n)=0\)
    \item $S(n,2)=2^{n-1}-1$
    \item $S(n,k) \cdot k!$ = number of ways to color $n$ nodes using colors from $\displaystyle 1$ to $\displaystyle \displaystyle k$ such that each color is used at least once.
    \item An $r$-associated Stirling number of the second kind is the number of ways to partition a set of $n$ objects into $\displaystyle \displaystyle k$ subsets, with each subset containing at least $r$ elements. It is denoted by $S_r( n , k )$ and obeys the recurrence relation.
    $\displaystyle \displaystyle S_r(n+1,k) = k S_r(n,k) + \binom{n}{r-1} S_r(n-r+1,k-1)$
    \item Denote the n objects to partition by the integers $\displaystyle 1, 2, \dots, n$. Define the reduced Stirling numbers of the second kind, denoted $S^d(n, k)$, to be the number of ways to partition the integers $\displaystyle 1, 2, \dots, n$ into k nonempty subsets such that all elements in each subset have pairwise distance at least d.
    That is, for any integers i and j in a given subset, it is required that $|i - j| \geq d$. It has been shown that these numbers satisfy,
    \(S^d(n, k) = S(n - d + 1, k - d + 1), n \geq k \geq d\)
\end{enumerate}

\subsection{Bell number}

\begin{enumerate}[resume]
    \item Counts the number of partitions of a set.
    \item $B_{n+1}=\displaystyle\sum_{k=0}^{n}\left(\frac{n}{k}\right) \cdot B_{k}$
    \item $B_{n}=\displaystyle\sum_{k=0}^{n}S(n,k)$ ,where $S(n,k)$ is stirling number of second kind.
\end{enumerate}

\subsection{Math}

\begin{enumerate}[resume]
    \item $\displaystyle ab \mod\ ac=a(b \mod\ c)$
    \item $\displaystyle \sum_{i = 0}^n{i\cdot i!=(n + 1)! - 1}$.
    \item $\displaystyle a^k - b^k = (a - b) \cdot (a^{k - 1}b^0 + a^{k - 2}b^1 + \dots + a^0b^{k - 1})$
    \item $\displaystyle \min(a + b, c) = a + \min(b, c - a)$
    \item $|a - b| + |b - c| + |c - a| = 2 (\max (a, b, c) - \min (a, b, c))$
    \item $\displaystyle a \cdot b \leq c \rightarrow a \leq \left \lfloor \dfrac{c}{b} \right \rfloor$ is correct
    \item $\displaystyle a \cdot b < c \rightarrow a <  \left \lfloor \dfrac{c}{b} \right \rfloor$ is incorrect
    \item $\displaystyle a \cdot b \geq c \rightarrow a \geq  \left \lfloor \dfrac{c}{b} \right \rfloor$ is correct
    \item $\displaystyle a \cdot b > c \rightarrow a >  \left \lfloor \dfrac{c}{b} \right \rfloor$ is correct
    \item For positive integer $n$, and arbitrary real numbers $m,x$,
    $\displaystyle \left \lfloor \dfrac{\left \lfloor x/m \right \rfloor}{n} \right \rfloor = \left \lfloor \dfrac{x}{mn} \right \rfloor$
    $\displaystyle \left \lceil \dfrac{\left \lceil x/m \right \rceil}{n} \right \rceil = \left \lceil \dfrac{x}{mn} \right \rceil$
    \item Lagrange’s identity:
    \[\displaystyle \displaystyle {\displaystyle {\begin{array}{@{}l}\left(\sum \limits_{k=1}^{n}a_{k}^{2}\right)\left(\sum \limits_{k=1}^{n}b_{k}^{2}\right)-\left(\sum\limits_{k=1}^{n}a_{k}b_{k}\right)^{2} \\ =\sum \limits_{i=1}^{n-1}\sum \limits_{j=i+1}^{n}\left(a_{i}b_{j}-a_{j}b_{i}\right)^{2} \\ ={\frac {1}{2}}\sum \limits_{i=1}^{n}\sum \limits_{j=1,j\neq i}^{n}(a_{i}b_{j}-a_{j}b_{i})^{2}\end{array}}}\]
    \item $\displaystyle \sum_{i = 1}^n{ia^i} = \frac{a(n a^{n + 1} - (n + 1) a^n + 1)}{(a - 1)^2}$
    \item \textbf{Vieta’s formulas:}
    Any general polynomial of degree $n$
    \[\displaystyle p(x) = a_nx^n + a_{n - 1}x^{n - 1} + \dots + a_1x + a_0\]
    (with the coefficients being real or complex numbers and $\displaystyle a_n \neq 0$) is known by the fundamental theorem of algebra to have $n$ (not necessarily distinct) complex roots $r_1, r_2,\dots, r_n$.
    \[\begin{cases}
        \text{$r_1 + r_2 + \dots + r_{n - 1} + r_n = - \frac{a_{n - 1}}{a_n}$}
        \\
\begin{array}{@{}l}
(r_1r_2 + r_1r_3 + \dots + r_1r_n) \\
+ (r_2r_3 + r_2r_4 + \dots + r_2r_n) + \dots \\
+ r_{n-1}r_n
= \dfrac{a_{n-2}}{a_n}
\end{array}
        \\\vdots\\
        \text{$r_1r_2\dots r_n = (-1)^n\frac{a_0}{a_n}.$}
        \end{cases}\]
    Vieta’s formulas can equivalently be written as
    \[\displaystyle \sum_{1 \leq i_1 < i_2 < \dots <i_k \leq n} \Bigg(\prod_{j = 1}^k{r_{i_j}}\Bigg) = (-1)^k\frac{a_{n - k}}{a_n},\]
    \item We are given n numbers $\displaystyle a_1,a_2,\dots,a_n$ and our task is to find a value $\displaystyle x$ that minimizes the sum,
    \[|a_1 - x| + |a_2 - x| + \dots + |a_n - x|\]
    optimal $\displaystyle x=$ median of the array.
    if $n$ is even $\displaystyle x =$ $[$left median,right median$]$ i.e. every number in this range will work.
    For minimizing
    \[\displaystyle (a_1 - x)^2 + (a_2 - x)^2 + \dots + (a_n - x)^2\]
    optimal $\displaystyle x = \dfrac{(a_1 + a_2 + \dots + a_n)}{ n}$
    \item Given an array a of n non-negative integers. The task is to find the sum of the product of elements of all the possible subsets. It is equal to the product of $\displaystyle (a_i + 1)$ for all $\displaystyle a_i$
    \item \textbf{Pentagonal number theorem:}
    In mathematics, the pentagonal number theorem states that
    \[
\begin{array}{@{}l}
\prod_{n=1}^{\infty} (1 - x^n)
= \sum_{k=-\infty}^{\infty} (-1)^k x^{\frac{k(3k-1)}{2}} \\
= 1 + \sum_{k=1}^{\infty} (-1)^k
\left(
x^{\frac{k(3k+1)}{2}} + x^{\frac{k(3k-1)}{2}}
\right).
\end{array}
\]
    In other words,
\[
\begin{array}{@{}l}
(1 - x)(1 - x^2)(1 - x^3)\cdots
\\ = 1 - x - x^2 + x^5 + x^7 - x^{12} \\ - x^{15} + x^{22} + x^{26} - \cdots .
\end{array}
\]
    The exponents $\displaystyle 1, 2, 5, 7, 12,\cdots$ on the right hand side are given by the formula $g_k = \frac{k(3k - 1)}{2}$ for $\displaystyle \displaystyle k = 1, -1, 2, -2, 3, \cdots$ and are called (generalized) pentagonal numbers.
    It is useful to find the partition number in $O(n \sqrt{n})$
\end{enumerate}

\subsection{Fibonacci Number}

\begin{enumerate}[resume]
    \item $\displaystyle F_0=0, F_1=1$ and $\displaystyle F_n=F_{n-1}+F_{n-2}$
    \item $\displaystyle F_n=\sum\limits_{k=0}^{\lfloor\frac{n-1}{2}\rfloor}\binom{n-k-1}{k}$
    \item $\displaystyle F_n=\frac{1}{\sqrt{5}}(\frac{1+\sqrt{5}}{2})^n-\frac{1}{\sqrt{5}}(\frac{1-\sqrt{5}}{2})^n$
    \item $\displaystyle \sum\limits_{i=1}^{n}F_i=F_{n+2}-1$
    \item $\displaystyle \sum\limits_{i=0}^{n-1}F_{2i+1}=F_{2n}$
    \item $\displaystyle \sum\limits_{i=1}^{n}F_{2i}=F_{2n+1}-1$
    \item $\displaystyle \sum\limits_{i=1}^{n}F_{i}^{2}=F_nF_{n+1}$
    \item $\displaystyle F_mF_{n+1}-F_{m-1}F_{n}=(-1)^nF_{m-n}$
    $\displaystyle F_{2n}=F_{n+1}^2-F_{n-1}^2=F_n(F_{n+1}+F_{n-1})$
    \item $\displaystyle F_mF_n+F_{m-1}F_{n-1}=F_{m+n-1}$
    $\displaystyle F_mF_{n+1}+F_{m-1}F_{n}=F_{m+n}$
    \item A number is Fibonacci if and only if one or both of $\displaystyle (5 \cdot n^2 + 4)$ or $\displaystyle (5 \cdot n^2 - 4)$ is a perfect square
    \item Every third number of the sequence is even and more generally, every $\displaystyle \displaystyle k^{th}$ number of the sequence is a multiple of $\displaystyle F_k$
    \item $gcd(F_m, F_n)=F_{gcd(m, n)}$
    \item Any three consecutive Fibonacci numbers are pairwise coprime, which means that, for every n, $gcd(F_n, F_{n+1})=gcd(F_n, F_{n+2}), gcd(F_{n+1}, F_{n+2})=1$
    \item If the members of the Fibonacci sequence are taken $mod$ $n$, the resulting sequence is periodic with period at most $6n$.
\end{enumerate}

\subsection{Pythagorean Triples}

\begin{enumerate}[resume]
    \item A Pythagorean triple consists of three positive integers $\displaystyle a, b,$ and $\displaystyle C$, such that $\displaystyle a^{2}+b^{2}=c^{2}$. Such a triple is commonly written $\displaystyle (a, b, c)$
    \item Euclid’s formula is a fundamental formula for generating Pythagorean triples given an arbitrary pair of integers m and n with $m >n >0$. The formula states that the integers
    \[\displaystyle a = m^{2}-n^{2}, b = 2mn,  c = m^{2}+n^{2}\]
    form a Pythagorean triple. The triple generated by Euclid’s formula is primitive if and only if m and n are coprime and not both odd. When both m and n are odd, then a, b, and c will be even, and the triple will not be primitive; however, dividing a, b, and c by 2 will yield a primitive triple when m and n are coprime and both odd.
    \item The following will generate all Pythagorean triples uniquely:
  \[
\begin{array}{@{}l}
a = k\,(m^{2} - n^{2}), \\
b = k\,(2mn),
c = k\,(m^{2} + n^{2}).
\end{array}
\]
    where m, n, and k are positive integers with $m > n$, and with m and n coprime and not both odd.
    \item \textbf{Theorem:}
    The number of Pythagorean triples
    {a,b,n} with $max{a,b,n} = n$ is given by
    \[\displaystyle \dfrac{1}{2}\left( \displaystyle\prod _{ p^{\alpha} || n}\left( 2\alpha +1\right) -1\right)\]
    where the product is over all prime divisors p of the form $4k+1$.
    The notation $\displaystyle p^{\alpha} || n$ stands for the highest exponent $\displaystyle \alpha$ for which $\displaystyle p^{\alpha}$ divides $n$
    \textbf{Example:}
    For $n = 2\cdot3^{2}\cdot5^{3}\cdot7^{4}\cdot11^{5}\cdot13^{6},$ the number of Pythagorean triples with hypotenuse n is $\displaystyle \dfrac{1}{2}\left( 7.13-1\right) =45$.
    To obtain a formula for the number of Pythagorean triples with hypotenuse less than a specific positive integer N, we may add the numbers corresponding to each $n < N$ given by the Theorem. There is no simple way to compute this as a function of N.
\end{enumerate}

\subsection{Sum of Squares Function}

\begin{enumerate}[resume]
    \item The function is defined as
    \[
r_k(n)
=
\left|
\begin{array}{@{}l}
(a_1,a_2,\dots,a_k) \in \mathbf{Z}^k : \\
n = a_1^2 + a_2^2 + \dots + a_k^2
\end{array}
\right|
\]
    \item The number of ways to write a natural number as sum of two squares is given by $r_2(n)$. It is given explicitly by
    $r_{2}\left(n\right) = 4\left(d_{1}\left(n\right) - d_{3}\left(n\right)\right)$
    where d1(n) is the number of divisors of n which are congruent with 1 modulo 4 and d3(n) is the number of divisors of n which are congruent with 3 modulo 4.
    The prime factorization $n = 2^{g}p_{1}^{f_{1}}p_{2}^{f_{2}}. . .q_{1}^{h_{1}}q_{2}^{h_{2}}. . . ,$ where $\displaystyle p_{i}$ are the prime factors of the form $\displaystyle p_{i} \equiv 1$ (mod 4), and $q_{i}$ are the prime factors of the form $q_{i} \equiv 3$ (mod 4) gives another formula $r_{2}\left(n\right) = 4\left(f_{1}+1\right)\left(f_{2}+1\right)\dots,$ if all exponents $h_{1}$,$h_{2}$,\dots are even. If one or more $h_{i}$ are odd, then $r_{2}\left(n\right) = 0.$
    \item The number of ways to represent n as the sum of four squares is eight times the sum of all its divisors which are not divisible by 4, i.e.
    $r_{4}\left( n \right) = 8 \displaystyle\sum _{d| n;4 \nmid d}d$
    $r_{8}\left(n\right) = 16 \displaystyle\sum _{d|n} \left(-1\right)^{n+d}d^{3}$
\end{enumerate}
\subsection{Number Theory}
\begin{enumerate}[resume]
    \item for $i > j$, $\displaystyle \gcd(i, j)$ $=$ $\displaystyle \gcd(i -j, j)$ $\displaystyle \leq (i -j)$
    \item $\displaystyle \sum_{x = 1}^n \left[ d | x^k \right] = \left \lfloor\dfrac{n}{\prod_{i = 0}{p_i^{\left \lceil \dfrac{e_i}{k}\right \rceil}}}\right \rfloor$,
    where $d = \prod_{i = 0}{p_i^{e_i}}$. Here, $[a | b]$ means if $\displaystyle a$ divides $b$ then it is $\displaystyle 1$, otherwise it is $0$.
    \item The number of lattice points on segment $\displaystyle (x_1,y_1)$ to $\displaystyle (x_2,y_2)$ is $\displaystyle \gcd(abs(x_1-x_2),abs(y_1-y_2)) + 1$
    \item $\displaystyle (n-1)! \mod  n = n -1$ if n is prime, 2 if $n = 4$, $0$ otherwise.
    \item A number has odd number of divisors if it is perfect square
    \item The sum of all divisors of a natural number n is odd if and only if $n=2^r\cdot k^2$ where $r$ is non-negative and $\displaystyle \displaystyle k$ is positive integer.
    \item Let $\displaystyle a$ and $b$ be coprime positive integers, and find integers $\displaystyle a{\prime}$ and $b{\prime}$ such that $\displaystyle aa{\prime} \equiv 1 \mod b$ and $bb{\prime} \equiv 1 \mod a$. Then the number of representations of a positive integers (n) as a non negative linear combination of $\displaystyle a$ and $b$ is
    \[\displaystyle \frac{n}{ab}-\Big\{\frac{b{\prime} n}{a}\Big\}-\Big\{\frac{a{\prime} n}{b}\Big\} + 1\]
    Here, $\displaystyle {x}$ denotes the fractional part of $\displaystyle x$.
\[
\begin{array}{@{}l}
\sum \limits_{i=1}^{a}\sum \limits_{j=1}^{b}\sum \limits_{k=1}^{c} d(i \cdot j \cdot k)
\\ =
\displaystyle \sum \limits_{\gcd (i,j)=\gcd (j,k) = \gcd (k,i) =1}
\displaystyle \left \lfloor \dfrac{a}{i}\right \rfloor \left \lfloor \dfrac{b}{j}\right \rfloor \left \lfloor \dfrac{c}{k} \right \rfloor
\end{array}
\]
    Here, $d(x) =$ number of divisors of $\displaystyle x$.
    \item \textbf{Gauss’s generalization of Wilson’s theorem:},
    Gauss proved that,
    \[\displaystyle \prod \limits_{k=1 \atop \gcd(k,m)=1}^{m}\!\!k\ \]
    \[\displaystyle \equiv {\begin{cases}-1{\pmod {m}}&{\text{if }}m=4,\;p^{\alpha },\;2p^{\alpha }\\\;\;\,1{\pmod {m}}&{\text{otherwise}}\end{cases}}\]
    where $\displaystyle p$ represents an odd prime and $\displaystyle \alpha$ a positive integer. The values of $m$ for which the product is $-1$ are precisely the ones where there is a primitive root modulo $m$.
\end{enumerate}
\subsection{Divisor Function}
\begin{enumerate}[resume]
    \item $\displaystyle \sigma_x(n)=\sum\limits_{d\vert n}^{}d^x$
    \item It is multiplicative i.e if $\displaystyle \gcd(a, b)=1\to \sigma_x(ab)=\sigma_x(a)\sigma_x(b)$.
    \item \[\displaystyle \sigma_x(n)=\prod\limits_{i=1}^{\tau}\frac{p_i^{(a_i+1)x}-1}{p_i^x-1}\]
    \item \textbf{Divisor Summatory Function}
    \begin{enumerate}
        \item Let $\displaystyle \sigma_0(k)$ be the number of divisors of $\displaystyle \displaystyle k$.
        \item $D(x)=\sum\limits_{n\leq{x}}{\sigma_0(n)}$
        \item $D(x)=\sum\limits_{k=1}^{x}\lfloor\frac{x}{k}\rfloor=2\sum\limits_{k=1}^{u}\lfloor\frac{x}{k}\rfloor-u^2$, where $u=\sqrt{x}$
        \item $D(n)=$Number of increasing arithmetic progressions where $n+1$ is the second or later term. (i.e. The last term, starting term can be any positive integer $\displaystyle \le n$. For example, $D(3)=5$ and there are 5 such arithmetic progressions: $\displaystyle (1, 2, 3, 4); (2, 3, 4); (1, 4); (2, 4); (3, 4).$
    \end{enumerate}
    \item Let $\displaystyle \sigma_1(k)$ be the sum of divisors of k. Then, $\displaystyle \sum\limits_{k=1}^{n}\sigma_1(k)=\sum\limits_{k=1}^{n}{k \left\lfloor \frac{n}{k} \right\rfloor}$
    \item $\displaystyle \prod\limits_{d\vert n}^{}d={n^{\frac{\sigma_0}{2}}}$ if $n$ is not a perfect square, and $=\sqrt{n} \cdot n^{\frac{\sigma_0-1}{2}}$ if $n $ is a perfect square.
\end{enumerate}
\subsection{Euler’s Totient function}
\begin{enumerate}[resume]
    \item The function is multiplicative.
    This means that if $\displaystyle \gcd(m, n) = 1$, $\displaystyle \phi(m \cdot n) = \phi(m) \cdot \phi(n)$.
    \item $\displaystyle \phi(n) = n\prod_{p |n}(1 - \frac{1}{p} )$
    \item If $p$ is prime and $(k \geq 1)$, then:
    $\displaystyle \phi(p^k) = p^{k - 1}(p - 1) = p^k(1 - \frac{1}{p})$
    \item $J_k(n)$, the Jordan totient function, is the number of $\displaystyle \displaystyle k$-tuples of positive integers all less than or equal to n that form a coprime $\displaystyle (k + 1)$-tuple together with $n$. It is a generalization of Euler’s totient, $\displaystyle \phi(n) = J_1(n)$.
    \(J_k(n) = n^k\prod_{p |n}(1 - \frac{1}{p^k})\)
    \item $\displaystyle \sum_{d|n}J_k(d) = n^k$
    \item $\displaystyle \sum_{d|n}\phi(d) = n$
    \item $\displaystyle \phi(n) = \sum_{d|n}\mu(d)\cdot\frac{n}{d} = n\sum_{d|n}\frac{\mu(d)}{d}$
    \item $\displaystyle \phi(n) = \sum_{d|n}d\cdot\mu(\frac{n}{d})$
    \item $\displaystyle \displaystyle {a}\vert{b} \to \varphi(a)\vert{\varphi(b)}$
    \item $n\vert{\varphi(a^n - 1)}$ for $\displaystyle a, n > 1$
    \item $\displaystyle \varphi(mn)=\varphi(m)\varphi(n)\cdot\frac{d}{\varphi(d)}$ where $d=gcd(m, n)$
    Note the special cases
    \[\varphi(2m)=
    \begin{cases}
      2\varphi(m) & ; if \, m  \, is \, even\\
      \varphi(m) & ; if \, m \, is \, odd\\
    \end{cases}\]
    \[\displaystyle \varphi(n^m)=n^{m-1}\varphi(n)\]
    \item $\displaystyle \varphi(lcm(m, n)) \cdot \varphi(gcd(m, n))=\varphi(m) \cdot \varphi(n)$
    Compare this to the formula $lcm(m, n)\cdot gcd(m,n)=m\cdot n$
    \item $\displaystyle \varphi(n)$ is even for $n\geq3$. Moreover, if if $n$ has $r$ distinct odd prime factors, $\displaystyle 2^r \vert \varphi(n)$
    \item $\displaystyle \sum\limits_{d\vert{n}}^{} \frac{\mu^2(d)}{\varphi(d)}=\frac{n}{\varphi(n)}$
    \item $\displaystyle \sum\limits_{1\leq k \leq n, \gcd(k, n)=1}k=\frac{1}{2}n\varphi(n)$ for $n > 1$
    \item $\displaystyle \frac{\varphi(n)}{n}=\frac{\varphi(rad(n))}{rad(n)}$ where $rad(n)=$
    $\displaystyle \prod\limits_{p\vert n, p \, prime} p$
    \item $\displaystyle \phi(m) \geq \log_2{m}$
    \item $\displaystyle \phi(\phi(m))\leq\frac{m}{2}$
    \item When $\displaystyle x \geq \log_2m$, then
    \[n^x\mod m=n^{\phi(m) + x\mod \phi(m)}\mod m\]
    \item $\displaystyle \sum\limits_{1\leq k\leq n, \gcd(k, n)=1}\gcd(k-1, n)=\varphi(n)d(n)$ where $d(n)$ is number of divisors. Same equation for $\displaystyle \gcd(a \cdot k-1, n)$ where $\displaystyle a$ and $n$ are coprime.
    \item For every $n$ there is at least one other integer $m\ne n$ such that $\displaystyle \varphi(m)=\varphi(n).$
    \item $\displaystyle \sum\limits_{i=1}^{n} {\varphi(i) \cdot \lfloor\frac{n}{i}\rfloor=\frac{n*(n+1)}{2}}$
    \item $\displaystyle \sum\limits_{i=1, i \% 2 \ne 0 }^{n} \varphi(i) \cdot \lfloor\frac{n}{i}\rfloor=\sum\limits_{k\geq 1}[\frac{n}{2^k}]^2$.
    Note that $[\, ]$ is used here to denote round operator not floor or ceil
    \item \[\displaystyle \sum\limits_{i=1}^{n}\sum\limits_{j=1}^{n}ij[\gcd(i, j)=1]=\sum\limits_{i=1}^{n}\varphi(i)i^2\]
    \item Average of coprimes of $n$ which are less than $n$ is $\displaystyle \dfrac{n}{2}$.
\end{enumerate}
\subsection{Mobius Function and Inversion}
\begin{enumerate}[resume]
    \item It is a multiplicative function.
    \item \[\sum_{d|n}\mu(d) =
    \begin{cases}
      1 & ; n=1\\
      0 & ; n > 0
    \end{cases}\]
    \item $\displaystyle \displaystyle \sum\limits_{n=1}^N {\mu}^2(n) = \displaystyle \sum\limits_{n=1}^{\sqrt{N}} \mu(k) \cdot \left \lfloor \dfrac{N}{k^2} \right\rfloor$
    This is also the number of square-free numbers $\displaystyle \le n$
    \item \textbf{Mobius inversion theorem:} The classic version states that if g and f are arithmetic functions satisfying
    $g(n) = \displaystyle \sum_{d|n}f(d) $ for every integer $n \ge 1$ then $g(n) = \displaystyle \sum_{d|n}\mu(d)g\left(\frac{n}{d}\right)$ for every integer $n \ge 1$
    \item If $\displaystyle F(n) = \displaystyle \prod_{d|n} f(d)$, then $\displaystyle F(n) = \displaystyle \prod_{d|n} F\left(\frac{n}{d}\right)^{\mu(d)}$
    \item $\displaystyle \displaystyle \sum_{d|n}\mu(d)\phi(d) = \displaystyle \prod_{j=1}^K (2 - P_j)$
    where $\displaystyle p_j$ is the primes factorization of $d$
    \item If $\displaystyle F(n)$ is multiplicative, $\displaystyle F \not\equiv 0$, then $\displaystyle \displaystyle \sum_{d|n} \mu(d) f(d) = \displaystyle \prod_{i=1} (1 - f(P_i)) \cdot$
    where $\displaystyle p_i$ are primes of $n$.
\end{enumerate}
\subsection{GCD and LCM}
\begin{enumerate}[resume]
    \item $\displaystyle \gcd(a, 0) = a$
    \item $\displaystyle \gcd(a, b) = \gcd(b, a \mod b)$
    \item Every common divisor of $\displaystyle a$ and $b$ is a divisor of $\displaystyle \gcd(a,b)$.
    \item if $m$ is any integer, then $\displaystyle \gcd(a + m {\cdot} b, b) = \gcd(a, b)$
    \item The gcd is a multiplicative function in the following sense: if $\displaystyle a_1$ and $\displaystyle a_2$ are relatively prime, then $\displaystyle \gcd(a_1 \cdot a_2, b) = \gcd(a_1, b) \cdot \gcd(a_2,b )$.
    \item $\displaystyle \gcd(a, b){\cdot} \operatorname{lcm}(a, b) = |a{\cdot}b|$
    \item $\displaystyle \gcd(a, \operatorname{lcm}(b, c)) = \operatorname{lcm}(\gcd(a, b), \gcd(a, c))$.
    \item $\displaystyle \operatorname{lcm}(a, \gcd(b, c)) = \gcd(\operatorname{lcm}(a, b), \operatorname{lcm}(a, c))$.
    \item For non-negative integers $\displaystyle a$ and $b$, where $\displaystyle a$ and $b$ are not both zero, $\displaystyle \gcd({n^a} - 1, {n^b} - 1) = n^{\gcd(a,b)} - 1$
    \item $\displaystyle \gcd(a, b) = \displaystyle \sum_{k|a \, \text{and} \, k|b} {\phi(k)}$
    \item $\displaystyle \displaystyle \sum_{i=1}^n [\gcd(i, n) = k] = { \phi{\left(\frac{n}{k}\right)}}$
    \item $\displaystyle \displaystyle \sum_{k=1}^n \gcd(k, n) = \displaystyle \sum_{d|n} d \cdot {\phi{\left(\frac{n}{d}\right)}}$
    \item $\displaystyle \displaystyle \sum_{k=1}^n x^{\gcd(k,n)} = \displaystyle \sum_{d|n} x^d \cdot {\phi{\left(\frac{n}{d}\right)}}$
    \item $\displaystyle \displaystyle \sum_{k=1}^n \frac{1}{\gcd(k, n)} = \displaystyle \sum_{d|n} \frac{1}{d} \cdot {\phi{\left(\frac{n}{d}\right)}} = \frac{1}{n} \displaystyle \sum_{d|n} d \cdot \phi(d)$
    \item $\displaystyle \displaystyle \sum_{k=1}^n \frac{k}{\gcd(k, n)} = \frac{n}{2} \cdot \displaystyle \sum_{d|n} \frac{1}{d} \cdot {\phi{\left(\frac{n}{d}\right)}} = \frac{n}{2} \cdot \frac{1}{n} \cdot \displaystyle \sum_{d|n} d \cdot \phi(d)$
    \item $\displaystyle \displaystyle \sum_{k=1}^n \frac{n}{\gcd(k, n)} = 2 * \displaystyle \sum_{k=1}^n \frac{k}{\gcd(k, n)} - 1$, for $n > 1$
    \item $\displaystyle \displaystyle \sum_{i=1}^n \sum_{j=1}^n [\gcd(i, j) = 1] = \displaystyle \sum_{d=1}^n \mu(d) \lfloor {\frac{n}{d} \rfloor}^2$
    \item $\displaystyle \displaystyle \sum_{i=1}^n \displaystyle\sum_{j=1}^n \gcd(i, j) = \displaystyle \sum_{d=1}^n \phi(d) \lfloor {\frac{n}{d} \rfloor}^2$
    \item $\displaystyle \sum_{i=1}^n \sum_{j=1}^n i \cdot j[\gcd(i, j) = 1] = \sum_{i=1}^n \phi(i)i^2$
    \item $\displaystyle F(n) = \displaystyle \sum_{i=1}^n \displaystyle \sum_{j=1}^n \operatorname{lcm}(i, j) = \displaystyle \sum_{l=1}^n {\left(\frac{\left( 1 + \lfloor{\frac{n}{l} \rfloor} \right) \left( \lfloor{\frac{n}{l} \rfloor} \right)} {2} \right)}^2 \displaystyle \sum_{d|l} \mu(d)ld$
    \item $\displaystyle \gcd(\operatorname{lcm}(a, b), \operatorname{lcm}(b, c), \operatorname{lcm}(a, c)) = \operatorname{lcm}(\gcd(a, b), \gcd(b, c), \gcd(a, c))$
    \item $\displaystyle \gcd(A_L, A_{L+1}, \dots, A_R) = \gcd(A_L, A_{L+1} - A_L, \dots, A_R - A_{R-1})$.
    \item Given n, If $SUM = LCM(1,n) + LCM(2,n) + \dots + LCM(n,n)$
    then SUM = $\displaystyle \dfrac{n}{2}( \displaystyle\sum _{ d| n}\left( \phi \left( d\right) \times d\right) +1$
  \end{enumerate}
  \subsection{Legendre Symbol}
  \begin{enumerate}[resume]
    \item $\displaystyle \left(\frac{a}{p}\right)\equiv a^{\frac{p-1}{2}} \pmod{p} \text{ and } \left(\frac{a}{p}\right)\in\{-1,0,1\}.$
    \item The Legendre symbol is periodic in its first (or top) argument: if $\displaystyle a \equiv b \pmod{p}$, then \\
    $\displaystyle \left(\frac{a}{p}\right)=\left(\frac{b}{p}\right).$
    \item The Legendre symbol is a completely multiplicative function of its top argument: \\
    $\displaystyle \left(\frac{ab}{p} \right) = \left(\frac{a}{p} \right) \left(\frac{b}{p} \right)$
    \item The Fibonacci numbers $\displaystyle 1,1,2,3,5,8,13\dots$ are defined by the recurrence $\displaystyle F_1 = F_2 = 1,F_{n+1} = F_n + F_{n-1}.$ If $\displaystyle p$ is a prime number then \\
    $\displaystyle F_{p-\left(\frac{p}{5}\right)} \equiv 0 {\pmod{p}}, \;\;F_p \equiv \left(\frac{p}{5}\right) {\pmod{p}}.$
    \item Continuing from previous point, \\
    $\displaystyle \left(\frac{p}{5}\right)= \text{infinite concatenation of the sequence}$ \\
    $\left(1,-1,-1,1,0\right) {\text{from }}p \geq 1$.
    \item If $n = k^2$ is perfect square then $\displaystyle \left(\frac{n}{p}\right)=1$ for every odd prime except $\displaystyle \left(\frac{n}{k}\right)=0$ if k is an odd prime.
    \end{enumerate}
  \subsection{Miscellaneous}
  \begin{enumerate}[resume]
    \item $\displaystyle \displaystyle a+b = a \oplus b +2(a \& b)$.
    \item $\displaystyle \displaystyle a + b = a \mid b + a \& b$
    \item $\displaystyle \displaystyle a \oplus b = a\mid b - a \& b$
    \item $\displaystyle \displaystyle k_{th}$ bit is set in $\displaystyle x$ iff $\displaystyle x \mod 2^{k - 1}\geq 2^k$. It comes handy when you need to look at the bits of the numbers which are pair sums or subset sums etc.
    \item $\displaystyle \displaystyle k_{th}$ bit is set in $\displaystyle x$ iff $\displaystyle x \mod 2^{k - 1} - x \mod 2^k \neq\ 0$ ($ = 2^k$ to be exact). It comes handy when you need to look at the bits of the numbers which are pair sums or subset sums etc.
    \item $n \mod 2^i = n \& (2^i - 1)$
    \item $\displaystyle 1\oplus 2 \oplus 3 \oplus \cdots \oplus (4k - 1) = 0$ for any $\displaystyle \displaystyle k \ge 0$
    \item \textbf{Erdos Gallai Theorem:}
    The degree sequence of an undirected graph is the non-increasing sequence of its vertex degrees
    A sequence of non-negative integers $d_1 \geq d_2 \geq \cdots \geq d_n$ can be represented as the degree sequence of finite simple graph on $n$ vertices if and only if $d_1 + d_2 + \cdots + d_n$ is even and
    \[\displaystyle \sum_{i = 1}^k {d_i \leq k(k - 1)} + \sum_{i = k + 1}^n \min(d_i, k)\]
    holds for every $\displaystyle \displaystyle k$ in $\displaystyle 1 \leq k \leq n$.
    \item $f_n = \sum_{i=0}^{n} \binom{n}{i} g_i \Rightarrow g_n = \sum_{i=0}^{n} (-1)^{n-i} \binom{n}{i} f_i$
\end{enumerate}
\subsection{Traingle}
\begin{itemize}[leftmargin=*]
    \item \textbf{Sine Rule:} $\frac{a}{\sin A} = \frac{b}{\sin B} = \frac{c}{\sin C} = 2R $
    \item \textbf{Cosine Rule:} $c^2 = a^2 + b^2 - 2ab \cos C$
    \item \textbf{Area:} $\Delta = \sqrt{s(s-a)(s-b)(s-c)} \\ = \frac{1}{2}bc \sin A$
    \item \textbf{Circumradius ($R$):} $R = \frac{abc}{4\Delta} = \frac{a}{2\sin A}$
    \item \textbf{Inradius ($r$):} $r = \frac{\Delta}{s}$
    \item \textbf{Ex-radii ($r_a, r_b, r_c$):} $r_a = \frac{\Delta}{s-a}$
    \item \textbf{Centroid ($G$):} $\left( \frac{\sum x}{3}, \frac{\sum y}{3} \right)$
    \item \textbf{Incenter ($I$):} $\left( \frac{ax_1+bx_2+cx_3}{a+b+c}, \frac{ay_1+by_2+cy_3}{a+b+c} \right)$
    \item \textbf{Excenter ($I_A$):} Replace $a$ with $-a$ in Incenter formula.
    \item \textbf{Euler Line:} $O, G, H$ are collinear; $HG = 2GO$.
    \item \textbf{Median ($m_a$):} $4m_a^2 = 2b^2 + 2c^2 - a^2$ (Apollonius)
    \item \textbf{Angle Bisector ($l_a$):} $l_a = \frac{2bc}{b+c}\cos(A/2)$
    \item \textbf{Stewart's Theorem:} Cevian(a line from A to a) length $d$ dividing $a$ into $m, n$:
    \[ b^2m + c^2n = a(d^2 + mn) \]
    \item \textbf{Centroid ($G$):} Intersection of \textit{medians} (vertex to midpoint).
    \item \textbf{Incenter ($I$):} Intersection of \textit{angle bisectors}.
    \item \textbf{Circumcenter ($O$):} Intersection of \textit{perpendicular bisectors} of the sides.
    \item \textbf{Orthocenter ($H$):} Intersection of \textit{altitudes} (vertex perpendicular to opposite side).
\end{itemize}
\subsection{Circle}
\begin{itemize}[leftmargin=*]
    \item \textbf{General Eq:} $x^2+y^2+2gx+2fy+c=0$. Center $(-g, -f)$, $r=\sqrt{g^2+f^2-c}$.
    \item \textbf{Tangent at $P(x_1, y_1)$:} $xx_1 + yy_1 + g(x+x_1) + f(y+y_1) + c = 0$
    \item \textbf{Slope of tangent from $P(x, y)$:} $y + f = m(x + g) \pm r\sqrt{1 + m^2}$
    \item \textbf{Tangency Condition ($y=mx+c$):} $c^2 = r^2(1+m^2)$
    \item \textbf{Chord Length:} $L = 2\sqrt{r^2 - d^2}$ ($d = $ dist to center)
    \item \textbf{Power of Point $P$:} $d^2 - r^2$ ($d = $ dist $P$ to center) (Useful for determining if a point is inside or outside)
    \item \textbf{Secant Theorem:} $PA \cdot PB = PC \cdot PD$ (AB and CD are 2 lines drawn from point P over the circle)
    \item \textbf{Sector Area:} $\frac{1}{2}r^2\theta$ ($\theta$ in rads)
    \item \textbf{Segment Area (Cut by chord):} $\frac{1}{2}r^2(\theta - \sin\theta)$
\end{itemize}
\subsection{Points \& 2D Line}
\begin{itemize}[leftmargin=*]
    \item \textbf{Dist. point $(x_1, y_1)$ to line $ax+by+c=0$:} 
    \[ d = \frac{|ax_1+by_1+c|}{\sqrt{a^2+b^2}} \]
    \item \textbf{Dist. between parallel lines:} $d = \frac{|c_1-c_2|}{\sqrt{a^2+b^2}}$
    \item \textbf{Shoelace Area:} $\frac{1}{2} | \sum (x_i y_{i+1} - x_{i+1} y_i) |$
    \item \textbf{foot of Perpendicular $(h,k)$}:
    \[ \frac{h-x_1}{a} = \frac{k-y_1}{b} = -\frac{ax_1+by_1+c}{a^2+b^2} \]
    \item \textbf{Reflection $(h,k)$:} Replace $- (...)$ with $-2 (...)$ above.
    \item \textbf{Rotation of Axes (CCW $\theta$):}
    $x = X\cos\theta - Y\sin\theta, \quad y = X\sin\theta + Y\cos\theta$
\end{itemize}
\subsection{Vector}
\begin{itemize}[leftmargin=*]
    \item \textbf{Direction Cosines:} $l=\cos\alpha, m=\cos\beta, n=\cos\gamma$. $l^2+m^2+n^2=1$.
    \item \textbf{Dot Product:} $\mathbf{a} \cdot \mathbf{b} = |\mathbf{a}||\mathbf{b}|\cos\theta$
    \item \textbf{Projection of $\mathbf{a}$ on $\mathbf{b}$:} $(\mathbf{a} \cdot \hat{b})\hat{b}$
    \item \textbf{Cross Product:} Area $\triangle = \frac{1}{2}|\mathbf{a} \times \mathbf{b}|$.
    \item \textbf{Scalar Triple (Vol):} $[\mathbf{a}\mathbf{b}\mathbf{c}] = \mathbf{a} \cdot (\mathbf{b} \times \mathbf{c})$. Tet Vol $= \frac{1}{6}|[\dots]|$.
    \item \textbf{Vector Triple:} $\mathbf{a} \times (\mathbf{b} \times \mathbf{c}) = (\mathbf{a}\cdot\mathbf{c})\mathbf{b} - (\mathbf{a}\cdot\mathbf{b})\mathbf{c}$
\end{itemize}
\subsection{3D Line}
Line 1: $\mathbf{r} = \mathbf{a}_1 + \lambda\mathbf{b}_1$. Line 2: $\mathbf{r} = \mathbf{a}_2 + \mu\mathbf{b}_2$.
\begin{itemize}[leftmargin=*]
    \item \textbf{Angle Between Lines:} $\cos \theta = \frac{|\mathbf{b}_1 \cdot \mathbf{b}_2|}{|\mathbf{b}_1||\mathbf{b}_2|}$
    \item \textbf{Condition for Intersection:} Scalar Triple Product $[\mathbf{a}_2-\mathbf{a}_1, \mathbf{b}_1, \mathbf{b}_2] = 0$.
    \item \textbf{Intersection Point:} Solve $\mathbf{a}_1 + \lambda\mathbf{b}_1 = \mathbf{a}_2 + \mu\mathbf{b}_2$ for $\lambda, \mu$.
    \item \textbf{Shortest Dist (Skew):} $D = \left| \frac{(\mathbf{a}_2 - \mathbf{a}_1) \cdot (\mathbf{b}_1 \times \mathbf{b}_2)}{|\mathbf{b}_1 \times \mathbf{b}_2|} \right|$
    \item \textbf{Dist Between Parallel Lines:} $D = \frac{|(\mathbf{a}_2 - \mathbf{a}_1) \times \mathbf{b}|}{|\mathbf{b}|}$
    \item \textbf{Dist Point $P$ to Line:} $D = \frac{|(\mathbf{p} - \mathbf{a}) \times \mathbf{b}|}{|\mathbf{b}|}$
    \item \textbf{Foot of Perp $F$ from $P$:} $\mathbf{f} = \mathbf{a} + \left( \frac{(\mathbf{p}-\mathbf{a})\cdot \mathbf{b}}{|\mathbf{b}|^2} \right)\mathbf{b}$
    \item \textbf{Image $P'$ in Line:} $\mathbf{p'} = 2\mathbf{f} - \mathbf{p}$.
\end{itemize}
\subsection{Plane}
Through $\mathbf{a}$, normal $\mathbf{n}$. Eq: $(\mathbf{r} - \mathbf{a}) \cdot \mathbf{n} = 0$.
\begin{itemize}[leftmargin=*]
    \item \textbf{Cartesian:} $Ax + By + Cz + D = 0$. Normal $(A, B, C)$.
    \item \textbf{Dist Point $P$ to Plane:} $D = \frac{|Ax_1 + By_1 + Cz_1 + D|}{\sqrt{A^2+B^2+C^2}}$
    \item \textbf{Angle Between Planes:} $\cos \theta = \frac{|\mathbf{n}_1 \cdot \mathbf{n}_2|}{|\mathbf{n}_1||\mathbf{n}_2|}$
    \item \textbf{Angle Between Line \& Plane:} $\sin \phi = \frac{|\mathbf{b} \cdot \mathbf{n}|}{|\mathbf{b}||\mathbf{n}|}$
    \item \textbf{Intersection (Line of two planes):} Direction $\mathbf{v} = \mathbf{n}_1 \times \mathbf{n}_2$.
    \item \textbf{Family of Planes:} $P_1 + \lambda P_2 = 0$ (Planes through intersection line).
    \item \textbf{Foot of Perp from $P$:} $\frac{x-x_1}{A} = \frac{y-y_1}{B} = \frac{z-z_1}{C} = \frac{-(Ax_1+By_1+Cz_1+D)}{A^2+B^2+C^2}$
    \item \textbf{Image of Point in Plane:} Use $-2$ instead of $-1$ above.
    \item \textbf{Plane Bisectors:} $\frac{A_1x+B_1y+C_1z+D_1}{\sqrt{A_1^2+B_1^2+C_1^2}} = \pm \frac{A_2x+B_2y+C_2z+D_2}{\sqrt{A_2^2+B_2^2+C_2^2}}$
\end{itemize}
\subsection{Sphere}
\begin{itemize}[leftmargin=*]
    \item \textbf{Standard:} $(x-u)^2 + (y-v)^2 + (z-w)^2 = r^2$.
    \item \textbf{Diameter Form:} $(x-x_1)(x-x_2) + (y-y_1)(y-y_2) + (z-z_1)(z-z_2) = 0$.
    \item \textbf{Tangent Plane at $P$:} $(\mathbf{r}-\mathbf{c}) \cdot (\mathbf{p}-\mathbf{c}) = r^2$.
    \item \textbf{Circle of Intersection:} Radius $R = \sqrt{r^2 - d^2}$.
    \item \textbf{Spherical Cosine Rule:} $\cos a = \cos b \cos c + \sin b \sin c \cos A$.
\end{itemize}
\subsection{Solid}
\begin{itemize}[leftmargin=*]
    \item \textbf{Cone:} $V=\frac{1}{3}\pi r^2 h$, $SA = \pi r(r+l)$
    \item \textbf{Sphere:} $V=\frac{4}{3}\pi r^3$, $SA = 4\pi r^2$
    \item \textbf{Spherical Cap:} $SA = 2\pi r h$. $V = \frac{1}{3}\pi h^2(3r-h)$.
    \item \textbf{Frustum (Cone):} $V = \frac{1}{3}\pi h (R^2 + r^2 + Rr)$. $LSA = \pi(R+r)l$.
    \item \textbf{Pyramid:} $V = \frac{1}{3} (\text{Base Area}) \times h$. $LSA = \frac{1}{2} (\text{Perimeter}) \times l$.
    \item \textbf{Regular Tetrahedron:} Side $a$. $H = a\sqrt{2/3}$. $V = \frac{a^3}{6\sqrt{2}}$. Face Angle $\cos\theta = 1/3$.
\end{itemize}
\subsection{Trigonometry}
\begin{itemize}[leftmargin=*]
    \item $\sin(A\pm B) = \sin A\cos B \pm \cos A\sin B$
    \item $\cos(A\pm B) = \cos A\cos B \mp \sin A\sin B$
    \item $\cos 2A = \cos^2 A - \sin^2 A = 2\cos^2 A - 1$
    \item $\tan^{-1}x \pm \tan^{-1}y = \tan^{-1}\left( \frac{x \pm y}{1 \mp xy} \right)$
    \textit{(Note: For $+$, if $xy > 1$, add $\pi$ to the result.)}
    \item $2\tan^{-1}x = \tan^{-1}\left(\frac{2x}{1-x^2}\right) = \sin^{-1}\left(\frac{2x}{1+x^2}\right) = \cos^{-1}\left(\frac{1-x^2}{1+x^2}\right)$
    \item $\tan^{-1}x + \cot^{-1}x = \frac{\pi}{2}$
\end{itemize}
